% Part: incompleteness
% Chapter: incompleteness-provability
% Section: first-incompleteness-thm

\documentclass[../../../include/open-logic-section]{subfiles}

\begin{document}

\olfileid{inc}{inp}{1in}

\olsection{Teorema Ora Jangkep Kapisan}

Saiki kita bisa njlentrehake bukti asli G\"odel kanggo teorema ora
jangkep kapisan. Upamana $\Th{T}$ teori sembarang kang bisa
diaksiomatisasi kanthi komputabel ing basa kang ngembangake basa
aritmetika, saengga $\Th{T}$ ngemot aksioma-aksioma $\Th{Q}$. Iki
tegese, mligine, $\Th{T}$ merepresentasikake fungsi lan relasi kang
komputabel.

Kita wis ngandharake yen kanthi pangodean formula lan bukti minangka
bilangan kang lumrah, relasi $\Prf[\Th{T}](x,y)$ iku komputabel, ing
ngendi $\Prf[\Th{T}](x,y)$ bener yen lan mung yen $x$ iku bilangan
G\"odel saka !!a{derivation} kanggo !!{formula} kang nduweni bilangan
G\"odel~$y$ ing~$\Th{T}$. Malah, G\"odel bisa nuduhake yen kanggo
teori tartamtu kang dipikirake G\"odel, relasi iki rekursif primitif
nganggo dhaptar 45 fungsi lan relasi ing makalahe. Relasi kaping 45,
$x B y$, persis $\Prf[\Th{T}](x,y)$ kanggo pilihan~$\Th{T}$
tartamtu mau. Elinga yen nalika G\"odel nggunakake tembung
``rekursif'' ing makalahe, saiki kita bakal nggunakake istilah
``rekursif primitif.''

Amarga $\Prf[\Th{T}](x,y)$ komputabel, relasi iki bisa
direpresentasikake ing $\Th{T}$. Kita nggunakake
$\OPrf[\Th{T}](x,y)$ kanggo nyebut formula kang merepresentasikake
relasi iki. Tetepna $\OProv[\Th{T}](y)$ minangka formula
$\lexists[x][\OPrf[\Th{T}](x,y)]$. Iki nggambarake relasi kaping 46,
$\fn{Bew}(y)$, ing dhaptare G\"odel. Kaya kang dicathet G\"odel, iki
relasi siji-sijine kang ``ora bisa ditegesake rekursif.'' Mbokmenawa
tegese mangkene: saka definisine durung cetha yen relasi iki
komputabel; lan pangembangan sabanjure pancen nuduhake yen ora.

Upamana $\Th{T}$ teori kang !!{axiomatizable} lan ngemot~$\Th{Q}$.
Mula $\Prf[\Th{T}](x, y)$ bisa diputusake, saengga bisa
direpresentasikake ing~$\Th{Q}$ dening
!!a{formula}~$\OPrf[\Th{T}](x, y)$. Tetepna $\OProv[\Th{T}](y)$
minangka formula kang wis diterangake ing ndhuwur. Miturut lema titik
tetep, ana formula $!G_\Th{T}$ saengga $\Th{Q}$ (lan mula
$\Th{T}$) !!{derive}s
\begin{equation}
\ollabel{eqn:qpf}
!G_\Th{T} \liff \lnot \OProv[\Th{T}](\gn{!G_\Th{T}}).
\end{equation}
Gatekna yen $!G_\Th{T}$ pokoke nyatakake ``$!G_\Th{T}$
ora !!{derivable} ing~$\Th{T}$.'' 

\begin{lem}\ollabel{lem:cons-G-unprov}
Yen $\Th{T}$ iku teori konsisten kang !!{axiomatizable} lan
ngembangake~$\Th{Q}$, mula $\Th{T} \Proves/ !G_\Th{T}$.
\end{lem}

\begin{proof}
Upamana $\Th{T}$ !!{derive}s $!G_\Th{T}$. Mula pancen \emph{ana}
!!a{derivation}, saengga kanggo sawijining bilangan $m$, relasi
$\Prf[\Th{T}](m,\Gn{!G_\Th{T}})$ bener. Nanging banjur $\Th{Q}$
!!{derive}s ukara $\OPrf[\Th{T}](\num m, \gn{!G_\Th{T}})$. Mula
$\Th{Q}$ !!{derive}s
$\lexists[x][\OPrf[\Th{T}](x,\gn{!G_\Th{T}})]$, kang miturut definisi
iku $\OProv[\Th{T}](\gn{!G_\Th{T}})$. Miturut \olref{eqn:qpf},
$\Th{Q}$ !!{derive}s $\lnot !G_\Th{T}$, lan amarga $\Th{T}$
ngembangake $\Th{Q}$, mangkono uga~$\Th{T}$. Kita wis nuduhake yen
$\Th{T}$ !!{derive}s $!G_\Th{T}$, mula teori iki uga !!{derive}s
$\lnot !G_\Th{T}$, saengga teori iki ora konsisten.
\end{proof}

\begin{defn}
\ollabel{thm:oconsis-q}
Teori $\Th{T}$ diarani \emph{$\omega$-konsisten} yen kahanan iki
kelakon: yen $\lexists[x][!A(x)]$ iku ukara sembarang lan $\Th{T}$
!!{derive}s $\lnot !A(\num 0)$, $\lnot !A(\num 1)$,
$\lnot !A(\num 2)$, \dots mula $\Th{T}$ ora mbuktekake
$\lexists[x][!A(x)]$.
\end{defn}

Gatekna yen saben teori kang $\omega$-konsisten uga konsisten. Iki
langsung saka kasunyatan yen $\Th{T}$ ora konsisten, mula
$\Th{T} \Proves !A$ kanggo saben~$!A$. Mligine, yen $\Th{T}$ ora
konsisten, teori iki !!{derive}s $\lnot !A(\num n)$ kanggo saben~$n$
lan uga !!{derive}s~$\lexists[x][!A(x)]$. Dadi, yen $\Th{T}$ ora
konsisten, teori iki ora $\omega$-konsisten. Miturut kontraposisi, yen
$\Th{T}$ $\omega$-konsisten, teori iki mesthi konsisten.

\begin{lem}\ollabel{lem:omega-cons-G-unref}
Yen $\Th{T}$ iku teori $\omega$-konsisten kang !!{axiomatizable}
lan ngembangake~$\Th{Q}$, mula $\Th{T} \Proves/ \lnot !G_\Th{T}$.
\end{lem}

\begin{proof}
Kita nuduhake yen $\Th{T}$ !!{derive}s $\lnot !G_\Th{T}$, mula
teori iku ora $\omega$-konsisten. Upamana $\Th{T}$ !!{derive}s
$\lnot !G_\Th{T}$. Yen $\Th{T}$ ora konsisten, teori iki
ora $\omega$-konsisten, lan bukti rampung. Yen ora, $\Th{T}$
konsisten, mula teori iki ora !!{derive} $!G_\Th{T}$ miturut
\olref{lem:cons-G-unprov}. Amarga ora ana !!{derivation} kanggo
$!G_\Th{T}$ ing $\Th{T}$, $\Th{Q}$ !!{derive}s
\[
\lnot \OPrf[\Th{T}](\num 0, \gn{!G_\Th{T}}), \lnot \OPrf[\Th{T}](\num 1,
\gn{!G_\Th{T}}), \lnot \OPrf[\Th{T}](\num 2, \gn{!G_\Th{T}}), \dots
\]
lan mangkono uga~$\Th{T}$. Kosok baline, miturut \olref{eqn:qpf},
$\lnot !G_\Th{T}$ ekuivalen karo
$\lexists[x][\OPrf[\Th{T}](x,\gn{!G_\Th{T}})]$. Dadi $\Th{T}$
ora $\omega$-konsisten.
\end{proof}

\begin{prob}
  Saben teori kang $\omega$-konsisten iku konsisten. Tuduhna yen
  kosok baline ora mesthi bener, yaiku ana teori kang konsisten
  nanging ora $\omega$-konsisten. Lakonana kanthi nuduhake yen
  $\Th{Q} \cup \{\lnot !G_\Th{Q}\}$ konsisten nanging ora
  $\omega$-konsisten.
\end{prob}

\begin{thm}
\ollabel{thm:first-incompleteness} Upamana $\Th{T}$ teori sembarang
kang $\omega$-konsisten, !!{axiomatizable}, lan
ngembangake~$\Th{Q}$. Mula $\Th{T}$ ora jangkep.
\end{thm}

\begin{proof}
  Yen $\Th{T}$ $\omega$-konsisten, teori iki konsisten, mula
  $\Th{T} \Proves/ !G_\Th{T}$ miturut \olref{lem:cons-G-unprov}.
  Miturut \olref{lem:omega-cons-G-unref},
  $\Th{T} \Proves/ \lnot !G_\Th{T}$. Iki tegese $\Th{T}$ ora
  jangkep, amarga teori iki ora !!{derive}s $!G_\Th{T}$ lan uga ora
  $\lnot !G_\Th{T}$.
\end{proof}

\end{document}
